\documentclass[12pt,a4paper]{article}
\usepackage[utf8]{inputenc}
\usepackage[T1]{fontenc}
\usepackage[french]{babel}
\usepackage{amsmath,amssymb,amsthm}
\usepackage{geometry}
\usepackage{hyperref}
\geometry{margin=2.5cm}

\newtheorem{theoreme}{Théorème}
\newtheorem{proposition}{Proposition}
\newtheorem{lemme}{Lemme}
\newtheorem{corollaire}{Corollaire}
\newtheorem{axiome}{Axiome}
\newtheorem{definition}{Définition}

\title{%
  \textbf{Chapitre 2}\\[0.5em]
  \Large Mécanique Harmonique du Chaos Discret\\[0.5em]
  \large Théorie Mathématique de l'Univers est au Carré\\[0.5em]
  \normalsize Philippe Thomas Savard, 2026
}
\author{Philippe Thomas Savard}
\date{2026}

\begin{document}
\maketitle
\tableofcontents
\newpage

\begin{abstract}
Ce chapitre développe la mécanique harmonique du chaos discret, second pilier de la
théorie de Savard. Nous montrons que les irrégularités apparentes dans la distribution
des nombres premiers obéissent à une mécanique harmonique sous-jacente qui transforme
le chaos apparent en un ordre géométrique profond. La théorie révèle que le «chaos
discret» des nombres premiers est en réalité une superposition d'harmoniques dont
les fréquences sont déterminées par la géométrie spectrale du Chapitre~1.
\end{abstract}

\section{Introduction: Du Chaos à l'Ordre}

À première vue, les nombres premiers semblent distribués de façon chaotique parmi
les entiers. Pourtant, comme le révèle Savard, cette distribution obéit à une
\emph{mécanique harmonique} précise. Le «chaos» n'est que la manifestation discrète
d'un système continu sous-jacent.

\begin{definition}[Chaos Discret]
On appelle \emph{chaos discret} l'irrégularité observée dans la séquence des écarts
entre nombres premiers consécutifs:
\[
  g_k = P_{k+1} - P_k, \quad k \geq 1.
\]
Cette séquence $(g_k)_{k \geq 1}$ constitue la manifestation discrète du chaos harmonique.
\end{definition}

\section{Harmoniques de la Distribution des Nombres Premiers}

\subsection{Décomposition Harmonique}

La fonction de comptage des nombres premiers $\pi(x)$ admet une décomposition en
harmoniques via la formule explicite de Riemann:
\[
  \pi(x) = \mathrm{li}(x) - \sum_{\rho} \mathrm{li}(x^\rho) - \ln 2 +
  \int_x^{\infty} \frac{dt}{t(t^2-1)\ln t},
\]
où la somme porte sur tous les zéros non triviaux $\rho$ de $\zeta(s)$.

Chaque terme $\mathrm{li}(x^\rho)$ représente une \emph{harmonique} de la distribution
des nombres premiers. La mécanique harmonique de Savard analyse comment ces harmoniques
interagissent pour produire le chaos discret observé.

\begin{theoreme}[Mécanique des Harmoniques]
Les harmoniques de la distribution des nombres premiers forment un système mécanique
dont l'état est entièrement déterminé par les suites $A$ et $B$ du Chapitre~1.
En particulier, l'amplitude de la $k$-ième harmonique est proportionnelle à $1/k$.
\end{theoreme}

\subsection{Fréquences Fondamentales}

\begin{definition}[Fréquence Spectrale]
La \emph{fréquence spectrale} de la $k$-ième harmonique est définie par:
\[
  \omega_k = \Im(\rho_k), \quad k = 1, 2, 3, \ldots,
\]
où $\rho_k$ est le $k$-ième zéro non trivial de $\zeta(s)$ (ordonné par partie
imaginaire croissante).
\end{definition}

Les premières fréquences spectrales sont approximativement:
$\omega_1 \approx 14.135$, $\omega_2 \approx 21.022$, $\omega_3 \approx 25.011$, \ldots

\section{Le Chaos Discret comme Manifestation Géométrique}

\begin{proposition}[Origine Géométrique du Chaos]
Les écarts $g_k = P_{k+1} - P_k$ entre nombres premiers consécutifs sont entièrement
déterminés par la géométrie spectrale des suites $A$ et $B$. Plus précisément:
\[
  g_k \equiv 0 \pmod{2}, \quad k \geq 2,
\]
et la distribution des écarts obéit à la loi de Cramér:
\[
  \Pr[g_k > t \cdot \ln^2 P_k] \approx e^{-t}.
\]
\end{proposition}

\section{Équations de la Mécanique Harmonique}

La mécanique harmonique du chaos discret est régie par le système d'équations suivant:

\begin{align}
  \frac{d^2 \phi_k}{dt^2} + \omega_k^2 \phi_k &= F_k(t), \label{eq:harmonic}\\
  \phi_k(0) &= a_k, \quad \frac{d\phi_k}{dt}(0) = b_k,
\end{align}

où:
\begin{itemize}
  \item $\phi_k(t)$ est l'amplitude de la $k$-ième harmonique,
  \item $\omega_k$ est la fréquence spectrale correspondante,
  \item $F_k(t)$ est la force de couplage entre harmoniques,
  \item $a_k, b_k$ sont les conditions initiales déterminées par les suites $A$ et $B$.
\end{itemize}

\begin{lemme}[Conservation de l'Énergie Harmonique]
Le système harmonique \eqref{eq:harmonic} conserve l'énergie totale:
\[
  E = \sum_{k=1}^{\infty} \left(\frac{1}{2}\left(\frac{d\phi_k}{dt}\right)^2 +
  \frac{\omega_k^2}{2}\phi_k^2\right) < \infty.
\]
\end{lemme}

\section{Relation avec la Géométrie Spectrale}

La connexion entre la mécanique harmonique (Chapitre~2) et la géométrie spectrale
(Chapitre~1) est fournie par le théorème suivant:

\begin{theoreme}[Unité de la Théorie]
Les harmoniques de la distribution des nombres premiers, définies par les zéros de
$\zeta(s)$, correspondent exactement aux oscillations géométriques du spectre des
suites $A$ et $B$ autour de la droite critique $\Re(s) = 1/2$.
En d'autres termes, la mécanique harmonique du chaos discret et la géométrie du
spectre des nombres premiers sont deux descriptions équivalentes du même phénomène.
\end{theoreme}

\section{Applications et Exemples}

\subsection{Prédiction des Écarts entre Nombres Premiers}

En utilisant les premières $N$ harmoniques du système, on peut prédire les écarts
$g_k$ avec une précision croissante en $N$:
\[
  g_k^{(N)} = \sum_{j=1}^{N} A_j \cos(\omega_j \ln P_k + \phi_j),
\]
où $A_j$ et $\phi_j$ sont l'amplitude et la phase de la $j$-ième harmonique.

\subsection{Exemple Numérique}

Pour les premiers écarts entre nombres premiers:
\begin{center}
\begin{tabular}{c|c|c|c}
  $k$ & $P_k$ & $P_{k+1}$ & $g_k$ \\\hline
  1 & 2 & 3 & 1 \\
  2 & 3 & 5 & 2 \\
  3 & 5 & 7 & 2 \\
  4 & 7 & 11 & 4 \\
  5 & 11 & 13 & 2 \\
\end{tabular}
\end{center}

La mécanique harmonique prédit que ces écarts reflètent la superposition des
harmoniques fondamentales du spectre, conformément à la géométrie spectrale.

\section{Conclusion du Chapitre 2}

La mécanique harmonique du chaos discret révèle que l'irrégularité apparente dans
la distribution des nombres premiers cache une structure harmonique précise, entièrement
cohérente avec la géométrie spectrale développée au Chapitre~1. Le «chaos» discret
est la manifestation visible d'un ordre géométrique profond, accessible par la
théorie de Savard.

\bigskip
\noindent\textbf{Auteur:} Philippe Thomas Savard\\
\noindent\textbf{Depuis:} 2016\\
\noindent\textbf{Dépôt:} Théorie Mathématique — L'Univers est au Carré, 2026

\end{document}
